\chapter{Bulk states, flux measurements, and defects}
\label{narr:defects}

A plaquette observable can serve several different roles: an exact conserved
sector label, a fluctuating measurement, or a diagnostic evaluated in an
approximate state. Distinguishing these roles is necessary before interpreting
a decrease of flux expectation as a density of thermal vortices or interpreting
a local low-energy mode as a physical non-Abelian excitation.

% EDITORIAL QUESTION: When is a Wilson plaquette value a sector label?
\section{An involution and its measurement probabilities}
For a physical Hermitian plaquette operator with $W_p^2=\Id$, define
\begin{equation}
 P_{p,\pm}=\frac{\Id\pm W_p}{2},\qquad
 \rho_{p,-}=\Tr(\rho_sP_{p,-})=\frac{1-\langle W_p\rangle}{2}.
 \label{eq:narr:flux_probability}
\end{equation}
This identity is exact in any physical state. At zero field, the pure Kitaev
Hamiltonian commutes with its plaquette operators, so simultaneous eigenspaces
can label conserved flux sectors~\cite{Kitaev2006}. A generic magnetic field
breaks this conservation. The probability in Eq.~\eqref{eq:narr:flux_probability}
still exists, but interpreting it as an equilibrium gas of independent visons
requires a further physical model.

In a Majorana calculation the ordered physical spin product must first be
mapped into the chosen convention. A Pfaffian of a covariance evaluates a
Gaussian expectation only after that map is established. A product of six
mean link values is another approximation, not an operator-level replacement
for the physical plaquette in all states. These alternatives may agree in a
special frozen sector and disagree once correlations or projection matter.

% EDITORIAL QUESTION: What must be held fixed when a vortex interaction is compared?
\section{Defect states and energy differences}
Let $E_{2v}(R;L)$ denote the energy of a declared two-defect construction and
$E_{\rm ref}(L)$ a reference computed with the same Hamiltonian and geometry.
The excess $E_{2v}(R;L)-E_{\rm ref}(L)$ includes a pair-creation cost. An interaction
normalized to vanish at a large reference separation is instead
\begin{equation}
 V(R;L)=E_{2v}(R;L)-E_{2v}(R_{\rm ref};L).
 \label{eq:narr:interaction}
\end{equation}
The two quantities should not share a label without stating the subtraction.
On a torus $R_{\rm ref}$ is finite and the defects interact with periodic images;
a thermodynamic interaction requires additional size control.

A frozen gauge-link flip, a constrained relaxed saddle, and an unconstrained
self-consistent solution define different state constructions. Relaxation may
change the local texture or even erase the intended defect. To interpret an
energy comparison one must record which degrees of freedom were fixed, which
were varied, how flux was checked, and whether physical parity was enforced.
A localized auxiliary eigenvector alone supplies none of these conditions.

% EDITORIAL QUESTION: What separates a localized low mode from a braiding claim?
\section{Localization and transport are separate tests}
A small positive matrix eigenvalue can indicate a candidate localized mode.
Its spatial weight, separation dependence, and isolation from other modes must
be inspected in the same geometry. Physical parity can exclude a nominal
single-fermion excitation. Distinguishing splitting from an actual zero mode
also requires a size or separation limit rather than a display tolerance.

Braiding adds a family of states and a transport prescription. A computed
holonomy depends on how the relevant subspace is identified, how its gauge is
matched along the path, and whether the path remains separated from unwanted
states. It cannot be inferred from a static localization plot. The inherited
enclosing-loop test currently returns about $0.5904$ radians and fails its
original near-$\pi$ assertion. This failure is reproduced on the original
source; its unchanged assertions are retained as an opt-in failing research check after
the PI retired the unsupported guarantee from required CI. It does
not disprove a general theorem about physical Kitaev vortices, but it prevents
this implementation from being offered as its numerical certification.
